\section{MVK Case Study: Sheet Tracking in Excited-State Molecular Fields}
\label{sec:results-mvk}

We evaluate the method on the methyl vinyl ketone (MVK) trajectory used in the
continuous-scatterplot and image-moment study of time-varying bivariate molecular
fields~\cite{cspImageMomentsMVK}.  That study reported a prominent transition
near the approach to the conical intersection, with dense sampling around
approximately \(29\)--\(31\,\mathrm{fs}\), and identified changes associated with
the carbonyl oxygen and the vinyl C3/C4 atoms.  Our goal is not to reproduce the
same atom-track analysis, but to test whether Reeb-space sheet tracking exposes
the same temporal transition as a feature-level event, and whether the domain
and range correspondence modes provide complementary information.

The MVK data are represented by two scalar fields, `orb00' and `orb01'.  For
each timestep, we compute the Reeb space, keep the top 20 sheets by area, and
track sheets between adjacent timesteps.  The viewer labels timesteps by their
simulation step and converts them to femtoseconds using
\[
  t_{\mathrm{fs}} = \frac{\mathrm{step}}{41.341374575751}.
\]
Thus the dense timestep window \(1240\)--\(1280\) corresponds to approximately
\(30.00\)--\(30.96\,\mathrm{fs}\).

\subsection{Range-Space Events Recover the Transition Window}

The clearest MVK result is obtained from the range-space metric.  In the current
configuration, the combined range score assigns all weight to sheet
intersection-over-union in the bivariate range.  A low best-continuation score
therefore indicates that a prominent sheet in one timestep does not have a
geometrically similar sheet in the next timestep.  The strongest range-event
scores occur near the transition interval reported in the prior
continuous-scatterplot analysis.

\begin{table}[t]
  \centering
  \caption{Important MVK intervals identified by Reeb-space sheet tracking.
  Range events are computed with the range IoU score and threshold
  \(\theta_R=0.5\).  Domain events use normalized domain-overlap vertex counts
  with the runtime median threshold \(\theta_D\).}
  \label{tab:mvk-events}
  \begin{tabular}{llrrrl}
    \toprule
    Stage & Mode & Interval & Time (fs) & Score & Observation \\
    \midrule
    MVK\_s1 & range & 1180--1200 & 28.54--29.03 & 36.66 &
    onset of strong range reconfiguration \\
    MVK\_s1 & range & 1200--1220 & 29.03--29.51 & 35.77 &
    strongest MVK\_s1 range continuation loss \\
    MVK\_s1 & range & 1276--1280 & 30.86--30.96 & 31.70 &
    late dense-window range reconfiguration \\
    MVK\_s2 & range & 1140--1160 & 27.58--28.06 & 37.21 &
    earliest strong MVK\_s2 range reconfiguration \\
    MVK\_s2 & range & 1160--1180 & 28.06--28.54 & 30.71 &
    continuation of the same event window \\
    MVK\_s2 & range & 1340--1360 & 32.41--32.90 & 22.54 &
    post-transition range reconfiguration \\
    \midrule
    MVK\_s1 & domain & 460--480 & 11.13--11.61 & 49.34 &
    early spatial-support churn \\
    MVK\_s1 & domain & 960--980 & 23.22--23.71 & 49.11 &
    mid-trajectory support change \\
    MVK\_s2 & domain & 320--340 & 7.74--8.22 & 52.63 &
    early support reorganization \\
    MVK\_s2 & domain & 300--320 & 7.26--7.74 & 52.60 &
    early support reorganization \\
    \bottomrule
  \end{tabular}
\end{table}

For MVK\_s1, the highest range-event scores are at
\(28.54\)--\(29.03\,\mathrm{fs}\), \(29.03\)--\(29.51\,\mathrm{fs}\), and
\(30.86\)--\(30.96\,\mathrm{fs}\).  For MVK\_s2, the strongest events begin
slightly earlier, at \(27.58\)--\(28.06\,\mathrm{fs}\) and
\(28.06\)--\(28.54\,\mathrm{fs}\), with an additional post-transition event at
\(32.41\)--\(32.90\,\mathrm{fs}\).  This agreement with the known transition
window suggests that range-space Reeb sheets capture a feature-level
reconfiguration of the bivariate molecular field.

\begin{figure*}[t]
  \centering
  \fbox{\parbox{0.92\linewidth}{
  Placeholder: range-mode Sankey views for MVK\_s1 and MVK\_s2 over
  \(27.5\)--\(31.5\,\mathrm{fs}\).  Highlight MVK\_s1 intervals
  1180--1200, 1200--1220, and 1276--1280, and MVK\_s2 intervals
  1140--1160 and 1160--1180.  The figure should show the range event graph
  above the Sankey panel, with selected intervals highlighted in the graph and
  in the corresponding columns of the Sankey diagram.}}
  \caption{Range-space Reeb-sheet events concentrate near the MVK transition
  window reported by prior continuous-scatterplot analysis.}
  \label{fig:mvk-range-events}
\end{figure*}

\subsection{Long-Lived Reeb-Sheet Feature Families}

The range metric also identifies persistent sheet families across the trajectory.
In MVK\_s1, three high-ranking sheets persist across the full time span
\(0.00\)--\(35.80\,\mathrm{fs}\): sheet \(0\rightarrow6341\) with mean
continuation score \(0.978\), sheet \(5251\rightarrow8\) with mean score
\(0.950\), and sheet \(4\rightarrow5\) with mean score \(0.896\).  In MVK\_s2,
sheet \(24\rightarrow503\) persists across the full trajectory with mean score
\(0.901\).  Two other prominent MVK\_s2 tracks are sheet \(22\rightarrow1892\),
which persists until \(27.58\,\mathrm{fs}\), and sheet \(2\rightarrow0\), which
starts at \(13.06\,\mathrm{fs}\) and persists to the end of the sequence.

These long-lived tracks are a central distinction from timestep-level
continuous-scatterplot summaries.  The method does not only say that the
scatterplot changes near a particular time; it identifies which prominent Reeb
sheets remain coherent over time and where their continuation becomes weak.

\begin{figure}[t]
  \centering
  \fbox{\parbox{0.88\linewidth}{
  Placeholder: selected continuing-feature graph and Sankey highlights.
  Show MVK\_s1 range tracks starting at sheets 0, 5251, and 4, and MVK\_s2
  range tracks starting at sheets 24, 22, and 2.}}
  \caption{Long-lived range-space Reeb-sheet tracks in MVK.}
  \label{fig:mvk-long-lived-tracks}
\end{figure}

\subsection{Domain-Space Tracking Answers a Different Question}

The domain mode uses overlap of regular mesh vertices between sheets.  This asks
whether the same spatial support continues to participate in corresponding
sheets.  It should not be expected to peak at exactly the same intervals as the
range metric, because a sheet can preserve its range-space footprint while its
supporting domain vertices change, or it can preserve spatial support while its
range-space geometry changes.

For MVK, the strongest domain-overlap event scores occur earlier than the main
range transition.  MVK\_s1 has high domain events at
\(11.13\)--\(11.61\,\mathrm{fs}\), \(18.87\)--\(19.35\,\mathrm{fs}\), and
\(23.22\)--\(23.71\,\mathrm{fs}\).  MVK\_s2 has its strongest domain events at
\(6.29\)--\(8.22\,\mathrm{fs}\).  The separate domain-change diagnostic, which
uses source retention, target inheritance, split/merge counts, and a churn term,
shows the same pattern: MVK\_s1 domain support reorganizes most strongly around
\(13.5\)--\(17.9\,\mathrm{fs}\), while MVK\_s2 has a strong early support
reorganization around \(6.8\)--\(8.2\,\mathrm{fs}\).

This behavior is useful if presented as complementarity rather than
contradiction.  Range-space events recover the known transition window, while
domain-space events expose when the spatial support of prominent sheets changes.
For the molecular interpretation, domain events should be inspected with the
linked sheet and fiber-surface views before assigning them to specific atoms or
functional groups.

\subsection{Domain/Range Complementarity}

The tool also computes a domain/range complementarity diagnostic by comparing,
for each source sheet, the target selected by the range metric against the
target selected by the domain metric.  The current implementation is useful for
finding examples but should be de-emphasized as a quantitative paper result.
The reason is that the range score is measured on \([0,1]\), while the domain
max-overlap value is currently used as a percentage on \([0,100]\) in the
complementarity score.  This makes absolute score magnitudes difficult to
interpret.

For analysis, we also evaluated a normalized variant:
\[
  C = \frac{1}{2}(L_R + L_D)\min(q_R,q_D),
\]
where \(q_R\) is the best range score, \(q_D\) is the best domain score after
dividing the domain max-overlap percentage by 100, \(L_R\) is the range-score
loss incurred by choosing the domain-selected target, and \(L_D\) is the
domain-score loss incurred by choosing the range-selected target.  This
normalized score highlights several transition-window examples:

\begin{table}[t]
  \centering
  \caption{Representative normalized domain/range complementarity examples.
  A low agreement fraction means that few source sheets choose the same target
  under range and domain criteria.}
  \label{tab:mvk-complementarity}
  \begin{tabular}{llrrrr}
    \toprule
    Stage & Interval & Time (fs) & Score & Source sheet & Range/Domain targets \\
    \midrule
    MVK\_s1 & 1240--1244 & 29.99--30.09 & 0.737 & 978 & 1095 / 78 \\
    MVK\_s1 & 1220--1240 & 29.51--29.99 & 0.732 & 133 & 3 / 53 \\
    MVK\_s2 & 1244--1248 & 30.09--30.19 & 0.801 & 144 & 128 / 117 \\
    MVK\_s2 & 1220--1240 & 29.51--29.99 & 0.798 & 96 & 99 / 299 \\
    MVK\_s2 & 1248--1252 & 30.19--30.28 & 0.755 & 128 & 120 / 109 \\
    \bottomrule
  \end{tabular}
\end{table}

These examples suggest that, during the transition window, range and domain
criteria may prefer different continuations for the same source sheet.  This is
consistent with the intended interpretation of the two metrics.  However, the
rolling Spearman correlation between range-event and domain-event scores is not
smooth enough to serve as a main result: the mean correlation is \(-0.094\) for
MVK\_s1 and \(-0.190\) for MVK\_s2, with large local variation.  We therefore
use complementarity as an exploratory diagnostic and rely on the range event
and continuing-feature views for the primary MVK evidence.

\subsection{Sensitivity to Removing Near-Origin Scalar Values}

We also tested whether removing regular vertices whose scalar pair lies near the
range-space origin improves the domain analysis.  The tested thresholds were
\(|\mathrm{orb00}|,|\mathrm{orb01}|\leq 0.001, 0.003, 0.006,\) and \(0.01\).
At the requested threshold \(0.003\), the filter removes on average \(4.04\%\)
of regular vertices in MVK\_s1 and \(7.24\%\) in MVK\_s2.  The range results are
unchanged, as expected, and the domain-event rankings remain qualitatively the
same.  Larger thresholds remove substantially more vertices but do not produce a
cleaner domain or complementarity signal.  We therefore keep the baseline MVK
results and treat the low-origin exclusion as a sensitivity check rather than a
main analysis setting.

\subsection{Interpretation}

The MVK case study supports three conclusions.  First, range-space Reeb-sheet
tracking recovers the known transition window near \(29\)--\(31\,\mathrm{fs}\)
as a feature-level reconfiguration event.  Second, the method provides
long-lived sheet tracks, giving a temporal feature graph rather than only a
sequence of timestep descriptors.  Third, domain-space overlap provides a
complementary view of spatial support.  In MVK, the strongest domain-support
changes occur earlier than the strongest range events, indicating that the two
metrics answer distinct questions.

The current results do not by themselves prove atom-specific claims about C3,
C4, or oxygen.  To make those claims, the selected sheets and links in
Tables~\ref{tab:mvk-events} and~\ref{tab:mvk-complementarity} should be shown
with the linked sheet and fiber-surface image views, overlaid on the molecular
structure.  If those views localize the selected sheets near the carbonyl oxygen
or vinyl C3/C4 region, then the Reeb-space results can be directly connected to
the atom-level interpretation from the previous continuous-scatterplot study.

\begin{figure*}[t]
  \centering
  \fbox{\parbox{0.92\linewidth}{
  Placeholder: linked sheet/fiber-surface detail views for selected
  complementarity examples.  Suggested examples: MVK\_s1 interval 1240--1244,
  source sheet 978 with range target 1095 and domain target 78; MVK\_s2
  interval 1244--1248, source sheet 144 with range target 128 and domain target
  117.  Capture side-by-side sheet images and side-by-side fiber-surface images
  from the clicked-link detail panel.}}
  \caption{Selected MVK sheets for inspecting range/domain disagreement near
  the transition window.}
  \label{fig:mvk-detail-views}
\end{figure*}

